\section{Introduction}

An agent that works for a long time on one objective will eventually get something wrong,
and what happens next is a trust decision. It has to say which earlier steps still stand,
which are contaminated, and where to resume. Wrong in one direction it redoes sound work.
Wrong in the other it builds on a spoiled result, which is worse and much harder to see.

The architectural answer under study is the execution graph. A scanner captures the
artifacts a generated pipeline produced and the functions that produced them, and that
record becomes the durable state the agent reasons over. The appeal is that a graph built
from execution cannot be talked out of what happened. Recent work treats provenance of this
kind as the substrate for agent trust \cite{wang2026traces} and as the basis for
recoverable execution from recorded checkpoints \cite{zhuang2026agentrewind}.

We take the architecture seriously enough to ask what it is worth, and the answer we find
is narrower and more useful than the one usually assumed. \textbf{The graph pays off
through computation over its structure, and not through presentation to a model.} Those are
different uses of the same artifact and they do not succeed together. A selector that walks
the recorded dataflow to the causal root of a set of flagged boundaries is exact where a
baseline without that relation is useless. Acceptance criteria attached to boundaries catch
faults a model reading the same evidence does not. But handing a model a graph-selected
region of context, which is the use the substrate is most often justified by, cannot be
distinguished from handing it a size-matched random region, and accumulating history costs
two orders of magnitude more context for no gain.

Answering the question at all required a measurement the system does not have. Trust labels
here are produced by one model reviewing another, so nothing in the loop holds an
independent answer; when a reviewer calls a boundary trusted, no artifact can contradict it.
There is good reason not to take that on faith, since judges prefer text they find familiar
\cite{wataoka2024selfpreference} and a judge's agreement with itself can be high while its
agreement with the truth is not \cite{norman2026reliability}, and a verifier checking only
what it can express admits confident false positives \cite{helff2026gaming}. We supply the missing answer
by injecting a fault into a named boundary through the same scope machinery a repair must
pass, so the poisoned boundary is confined by mechanism rather than by assertion, and what
it contaminates is read from the pipeline's own dataflow. We are not aware of prior work
scoring an agent's own boundary trust labels against faults injected into its own pipeline
with known ground truth.

The paper's second contribution is methodological and, we think, the more transferable one.
Five times during this work a measurement produced a confident number that was an artifact
of how we had built the experiment rather than a property of the system: a static-code arm
that scored 75\,\% because our faults were source mutations and collapsed to 0\,\% once
they moved into recorded data; two models that returned a fixed answer regardless of input
while appearing to score respectably; and a selection result we reported as significant on a
third of its eventual sample and had to retract. Each would have survived into a published
table. We report them alongside the results because a benchmark for verifying agents that
cannot catch its own defects has no standing to measure anyone else's.
Extended related work appears in Appendix~\ref{app:related}.
